\subsection{Auxiliary linear systems for nonlinear optics equation}\label{subs2.2}

The nonlinear optics ($N$-wave) equation \eqref{NW1} is the compatibility condition
of the systems \eqref{au}, where
\begin{gather} \label{NW4}
G(x,t,z)=\I z D-\zeta(x,t), \qquad F(x,t,z)=\I z \wh D-\wh \zeta(x,t); \nonumber\\ \zeta=[D, \vr], \qquad \wh \zeta =[\wh D, \vr];
\end{gather}
see \cite{ZaMa} for the case $N=3$ and \cite{AbH} for $N>3$. We shall need
some preliminary results on the Weyl theory of the auxiliary system $y_x=Gy$
from \cite[Chapter~4]{SaSaR} (see also~\cite{SaA17}). The normalized fundamental solution~$w$
of such a system is def\/ined by the formula
\begin{gather} \label{NW5}
w_x(x,z)=\big(\I z D-\zeta(x)\big)w(x,z), \qquad w(0,z)=I_m, \qquad \zeta =- \zeta^* .
\end{gather}
Here and later we assume that $D$ is a f\/ixed matrix satisfying \eqref{NW2}. We consider
system~\eqref{NW5} with locally bounded potentials~$\zeta$, that is, potentials satisfying (for each $l<\infty$)
the inequality
\begin{gather}
  \sup\limits_{0 < x < l}
   \bigl\|
    \zeta (x)
   \bigr\|
  < \infty.
 \label{NW6}
\end{gather}
\begin{Definition} \label{NWDn2}
 A generalized Weyl function $($GW-function$)$ of system
\eqref{NW5}, where $\zeta$ is locally bounded, is an $m \times m$ matrix function $\vp$ such that for some $M>0$ it is analytic in the domain
$\BC^{-}_M = \{ z \colon \Im(z) < - M \}$ and the
inequality
\begin{gather}
  \sup\limits_{x \leq l, \, \Im (z) < -M}
   \bigl\|
    w(x,z)
    \varphi (z)\exp \{-\I z x D\}
   \bigr\|
 < \infty
 \label{NW7}
\end{gather}
holds for each
$l < \infty$.
\end{Definition}

\begin{Remark} We note that the Weyl function of the system \eqref{NW5} is def\/ined $\big($for $\zeta$ bounded on $[0,\, \infty)\big)$
by the analog \eqref{NWd6} of the inequality \eqref{2.0} and by normalization conditions \eqref{NWd7}. If~$\zeta$ is bounded on $[0,\, \infty)$, this Weyl function
coincides with the normalized {\rm GW}-function.
$($See the discussion after formula~\eqref{NWd7} and Def\/inition~\ref{NWDn2} of the {\rm GW}-function above.$)$ The fact that the Weyl function satisf\/ies~\eqref{NW7} explains the term ``generalized Weyl function'' $($or ``{\rm GW}-function''$)$.
\end{Remark}

The inverse spectral problem (ISpP) for system
\eqref{NW5} is the problem of recovering (from a~GW-function $\varphi$)
a potential
$\zeta (x) = - \zeta^* (x)$
 such that~\eqref{NW7} is valid and the diagonal entries~$\zeta_{kk}$ of~$\zeta$ equal zero $($i.e., $\zeta_{kk} \equiv 0)$.

\begin{Notation}
The notation $\mf M$
stands for
an operator mapping the pair $D$ and~$\varphi $
into the corresponding potential~$\zeta$
$($i.e., ${\mf M} (D, \varphi) = \zeta)$. In other words, ${\mf M} (D, \varphi)$ stands for a solution of the ISpP.
\end{Notation}


\noindent The following theorem (i.e., \cite[Theorem 4.8]{SaSaR}) is valid.
\begin{Theorem}\label{NWTmUniq}
For any
matrix function
$\varphi (z)$ which is
analytic and bounded in
$\BC^{-}_M$
and has the
property
\begin{gather}
  \int_{- \infty}^\infty
   \bigl(
    \varphi (z) - I_m
   \bigr)^*
   \bigl(
    \varphi (z) - I_m
   \bigr)
  d \xi
  < \infty,
\qquad
  z = \xi + \I \eta,
  \qquad
  \eta <- M ,
 \label{NW8}
\end{gather}
there is at most one solution of the ISpP.
\end{Theorem}
Our next corollary for systems on $(0,l)$ is immediate from the proof (see \cite[pp.~108--109]{SaSaR}) of Theorem~\ref{NWTmUniq}.

\begin{Corollary} \label{CyUniq} Assume that \eqref{NW6} is valid and let relations \eqref{NW7} and \eqref{NW8} hold
for an analytic and bounded matrix function~$\vp$. Then, $\zeta(x)$ is uniquely defined on~$(0,l)$.
\end{Corollary}

\begin{Remark}\label{NWInvPr} Under somewhat stronger $($than in~\eqref{NW8}$)$
restrictions, the solution of the ISpP always exists. Namely, if for some matrix $\phi_0$ and some $M>0$
 we have
\begin{gather} \label{NW9}
\sup_{\Im(z)<-M}
   \big\| \,
   z ( \varphi (z) - I_m) \,
   \big\|
   < \infty, \quad \det \, \varphi (z) \not= 0 \quad {\mathrm{for}} \quad \Im(z) <-M,
     \\
   (\xi+\I\eta)\big( \varphi (\xi+\I\eta) - I_m - \phi_0 / (\xi+\I \eta) \big)
   \in L_{m \times m}^2 (- \infty, \infty)
   \quad \text{for all f\/ixed} \
   \eta<-M,
 \label{NW10}
\end{gather}
then ${\mf M} (D, \varphi)$ is constructed in \cite[Theorem~4.10]{SaSaR}.
\end{Remark}

The situation becomes simpler when
$\zeta(x)$ is uniformly bounded on $[0, \infty)$, that is,
\begin{gather}
  \sup\limits_{0 < x < \infty}
   \bigl\| \zeta (x) \bigr\|
 \leq M_0 .
 \label{NWd8}
\end{gather}
We recall that  a Weyl function of system \eqref{NW5} is introduced in another way
than a GW-function. Namely, a Weyl function
 is an analytic $m \times m$ matrix function~$\varphi (z)$,
satisfying for certain
$M > 0$
and
$r > 0$
and for all~$z$
from the domain
$\BC^{-}_M = \{ z \colon \Im(z) < - M \}$
the inequality
\begin{gather}
  \int_0^\infty
      \exp \{\I \overline z x D \}
     \varphi (z)^*
     w (x,z)^*
     w (x,z)
     \varphi (z)
     \exp
      \bigl\{{-}
        (\I z D + r I_m )x
      \bigr\}
  dx
 < \infty,
 \label{NWd6}
\end{gather}
and the normalization conditions on the entries $\varphi_{ij} (z)$:
\begin{gather}
   \varphi_{ij} (z) \equiv 1 \qquad \mbox{for} \quad i= j, \qquad
   \varphi_{ij} (z) \equiv 0 \qquad \mbox{for} \quad i > j.
 \label{NWd7}
\end{gather}
When \eqref{NWd8} holds,
a Weyl function of system \eqref{NW5} exists and is unique.
Moreover, for that case~$\vp$ is the unique GW-function (of system \eqref{NW5} with the given $\zeta$)
satisfying normalization conditions~\eqref{NWd7}. In order to construct this Weyl (and simultaneously GW-) function
we use matrices $j$ of the form~\eqref{1.3} for each $1\leq m_1 < m$, that is, we set
\begin{gather*}
   %\label{NW11}
J_k:= \left[
     \begin{matrix}
       I_k & 0   \\
       0 & -I_{m-k}
     \end{matrix}
    \right], \qquad   1\leq k < m.
\end{gather*}
Now, the Weyl function is constructed \cite[pp.~103--106]{SaSaR}
in the following way.


First, for each $ k $, we introduce a class of
$m \times (m-k)$ matrix functions $\cQ_k$, which are meromorphic in some semi-plane $\BC^{-}_{M_k}$ and satisfy the inequalities
\begin{gather}
\cQ_k (z)^*
  \cQ_k(z)
 > 0, \qquad
  \cQ_k (z)^*
  J_k
  \cQ_k (z)
\leq 0,
 \label{NW12}
\end{gather}
excluding, possibly, isolated points. These $\cQ_k$ are called nonsingular with property-$J_k$.
Assu\-ming $M_k>2M_0/(d_k-d_{k+1})$ and using \eqref{NW12}, one can show that the matrix function
\begin{gather}
  \psi_k (x,z)
= \begin{bmatrix}I_k &0 \end{bmatrix}
 w (x,z)^{-1}
 \cQ_k (z)
 \left(
  \begin{bmatrix}0 &I_{m-k} \end{bmatrix}
  w (x,z)^{-1}
  \cQ_k(z)
 \right)^{-1}
 \label{NW13}
\end{gather}
is well-def\/ined for $x \geq 0$, $z \in \BC^{-}_{M_k}$, and satisf\/ies the inequality
\begin{gather}
\begin{bmatrix} \psi_k (z)^* &I_{m-k} \end{bmatrix}
 J_k
 \left[
  \begin{matrix}
    \psi_k (z) \\
    I_{m-k}
  \end{matrix}
 \right]
\leq 0 , \qquad {\mathrm{i.e.,}} \qquad \psi_k (z)^*
  \psi_k (z)
\leq I_{m-k} .
 \label{NW14}
\end{gather}
The set of matrices $\psi_k(x,z)$ given by \eqref{NW13}, where $x$ and $z$ are f\/ixed and matrix functions $\cQ_k(z)$ are nonsingular with property-$J_k$, is denoted
by $\cN_k(x,z)$. These sets are embedded and have a~point limit, that is, similar to \eqref{2.3} and \eqref{2.3!} we have
 \begin{gather}
  \breve \psi_k(z) = \bigcap\limits_{x < \infty}
           {\cal N}_k (x,z), \qquad \breve \psi_k(z)=\lim_{x\to \infty} \psi_k(x,z),
\qquad
  z \in \BC_{M_k}^{-} .
 \label{NW15}
\end{gather}
In this way we recover the $(k+1)$th column of the Weyl function $\varphi$. More precisely, we have
\begin{gather}
  \left\{
   \varphi_{i, k+1}(z)
  \right\}_{i=1}^k
 = \breve \psi_k (z)\begin{bmatrix} 1\\ 0 \\ \ldots \\0
 \end{bmatrix}, \qquad \Im(z)<-M \label{NW16}
\end{gather}
for any $M>\max\limits_{1\leq k<m}\big(2M_0/(d_k-d_{k+1})\big)$. There is also an inverse transformation \cite[Remark~4.6]{SaSaR}, which expresses $\breve \psi_k$ via $\varphi$:
\begin{gather}
  \breve\psi_k^{\,} (z)
= \begin{bmatrix} I_k & 0\end{bmatrix}
 \varphi (z)
 \left[
  \begin{matrix}
    0  \\
    I_{m-k}
  \end{matrix}
 \right]
 \biggl(
 \begin{bmatrix} 0 & I_{m-k}\end{bmatrix}
  \varphi (z)
  \left[
  \begin{matrix}
    0  \\
    I_{m-k}
  \end{matrix}
  \right]
 \biggr)^{-1} .
 \label{NW17}
\end{gather}
Finally, we will need a representation of $w(x,z)$ on intervals $0\leq x\leq l$, $l<\infty$, \cite[equation~(4.37)]{SaSaR} :
\begin{gather}
  w(x,z) = \exp \{ \I z x D\}
    + \int_{d_m x}^{d_1 x}
           \exp \{\I zt\}
     N(x,t)
     dt, \qquad \sup_{x\leq l} \| N(x,t) \| < \infty.
 \label{NW18}
\end{gather}

